\documentclass{article}
\usepackage[paperwidth=180mm,paperheight=160mm,margin=16mm]{geometry}
\usepackage{amsmath,unicode-math,xcolor}
\setmathfont{XITS Math}
\usepackage{experimental/xe-output}
\pagestyle{empty}
\begin{document}
\noindent\textbf{1. Native single-line box}\par\medskip
\noindent\begingroup\fboxsep=0pt\fboxrule=.3pt
\fbox{$\displaystyle\left(a+b+c\right)^2_i$}\endgroup
\par\bigskip
\noindent\textbf{2. XeTeX output-translation prototype}\par\medskip
\noindent\textcolor{red}{Red: original single-line box dimensions.}\par
\noindent\textcolor{blue}{Blue: revised outer box dimensions.}\par
\noindent Black: the actual formula content, emitted once.
\[
 \XeFlatShowBoxes
 \begin{xeflatmatchedsplit}
  \left(a+b\\+c\right)^2_i
 \end{xeflatmatchedsplit}
\]
\noindent The content list is still a single horizontal list. Only its outer
width and depth are changed. Output transformations move the second range
down and left. There are \textbf{no separate TeX row boxes} in this prototype.
\par\bigskip
\noindent The log contains the full box trees before and after the dimension
change, under \texttt{MSXE BEFORE} and \texttt{MSXE AFTER}.
\end{document}
